Suppose now that $v\in L^2(\Omega)$ is such that
\begin{equation}\label{eq:adjoint-weak-orthogonality-to-elliptic-range}
\int_\Omega(Lu)v\,\mathrm{d}\bm{x}=0
\end{equation}
for every $u\in H^{2m}(\Omega)$ such that $B_ju=0$ on $\partial\Omega$ for $j=1,\ldots,m$. If we actually knew that $v\in H^{2m}(\Omega)$, then we could use Green\textquotesingle s formula to conclude that
\begin{equation}\label{eq:adjoint-weak-orthogonality-green-consequence}
\int_\Omega u(L^*v)\,\mathrm{d}\bm{x}=-\sum_{j=1}^m\int_{\partial\Omega}S_juC_jv\,\mathrm{d}S.
\end{equation}
Since the $B_ju$ and $S_ju$ can be chosen arbitrarily and independently, \eqref{eq:adjoint-weak-orthogonality-green-consequence} implies that $L^*v=0$ and $C_jv=0$ on the boundary. Even without the assumption that $v\in H^{2m}(\Omega)$, we find $L^*v=0$ in the sense of distributions by restricting $u$ to $\mathcal{D}(\Omega)$. What remains to be done is to show that any $v\in L^2(\Omega)$ satisfying \eqref{eq:adjoint-weak-orthogonality-to-elliptic-range} actually is in $H^{2m}(\Omega)$.

Let
\begin{equation}\label{eq:adjoint-homogeneous-boundary-nullspace}
N^*=\{v\in H^{2m}(\Omega)\mid L^*v=0,\ C_1v=\cdots=C_mv=0\text{ on }\partial\Omega\},
\end{equation}
let $(N^*)^\perp$ be the orthogonal complement of $N^*$ in $L^2(\Omega)$ and let $M^*=(N^*)^\perp\cap H^{2m}(\Omega)$. On $M^*$ we consider the quadratic form
\begin{equation}\label{eq:adjoint-auxiliary-quadratic-form}
[u,v]=\int_\Omega L^*uL^*v\,\mathrm{d}\bm{x}+\sum_{j=1}^m(C_ju,C_jv)_{2m-m_j-1/2},
\end{equation}
where $m_j$ is the order of $C_j$ and $(\cdot,\cdot)_s$ denotes the inner product in $H^s(\partial\Omega)$. We shall show that this quadratic form is coercive on the space $M^*$.

\begin{Lemma}
\label{lem:adjoint-auxiliary-form-complement-coercivity}
There exists a constant $C$ such that, for all $u\in M^*$, we have
\begin{equation}\label{eq:adjoint-auxiliary-form-complement-coercivity}
\|u\|_{2m}^2\leq C[u,u].
\end{equation}

\end{Lemma}

\begin{Proof}
Since the $C_j$ boundary operators satisfy the complementing condition for $L^*$, we can use the Agmon--Douglis--Nirenberg theorem to get
\begin{equation}\label{eq:adjoint-auxiliary-form-garding-estimate}
\|u\|_{2m}^2\leq C([u,u]+\|u\|_0^2)
\end{equation}
for any $u\in H^{2m}(\Omega)$.

Suppose now for the sake of contradiction that $u_n\in M^*$ and $[u_n,u_n]\to0$, whereas $\|u_n\|_{2m}=1$. Then a subsequence of the $u_n$ converges weakly in $H^{2m}(\Omega)$ and strongly in $L^2(\Omega)$; let $u$ be the limit. It now follows from \eqref{eq:adjoint-auxiliary-form-garding-estimate} that the subsequence actually converges to $u$ strongly in $H^{2m}(\Omega)$. Since $[u,u]=0$, we have $u\in N^*$, and since $u_n\in M^*$, we also have $u\in M^*$. Since $N^*$ and $M^*$ are orthogonal in $L^2(\Omega)$, this implies $u=0$. But that is a contradiction since $\|u\|_{2m}=1$.
\end{Proof}

The Lax--Milgram lemma now implies that the equation $[u,v]=(f,v)$ for all $v\in M^*$ has a unique solution $u\in M^*$. If $f\in(N^*)^\perp$, we actually have $[u,v]=(f,v)$ for all $v\in H^{2m}(\Omega)$. We can summarize this as follows.
